\documentclass[12pt,a4paper]{article}

\usepackage[utf8]{inputenc}
\usepackage[T1]{fontenc}
\usepackage[spanish]{babel}
\usepackage{datatool}
\usepackage{booktabs}
\usepackage[table]{xcolor}
\usepackage{array} % Necesario para >{\raggedright\arraybackslash}
\usepackage[margin=1.5cm]{geometry} % Márgenes ajustados
\usepackage{fancyhdr} % Paquete para encabezados fijos


\definecolor{ctabla}{RGB}{205, 225, 225}

% 1. CARGA DE ARCHIVOS
\DTLloaddb[autokeys=false, headers={COD,Descripcion,Monto,Porcentaje}]{presupuesto}{porcentajes_presupuesto_2025.csv}
\DTLloaddb[autokeys=false, headers={Macrodistrito,Distrito,Monto,Porcentaje}]{fac}{macros.csv}

% --- CONFIGURACIÓN DEL ENCABEZADO ---
\pagestyle{fancy}
\fancyhf{}
\fancyhead[L]{ % Logo a la izquierda
%    \includegraphics[width=3.5cm]{logo_alcaldia.png} % Asegúrate de tener este archivo
}
\fancyhead[R]{ % Texto institucional a la derecha
    \small \textbf{FISCALIZACIÓN GAMLP} \\
    \small \textit{Gestión 2025}
}


\begin{document}
\section*{Presupuesto Municipal Gestión GAMLP 2025}
\begin{small}
\renewcommand{\arraystretch}{0.8} % Reduce el espacio entre filas para que quepa en una página
\centering
\begin{tabular}{l >{\raggedright\arraybackslash}p{9cm} r r} 
  \toprule
\textbf{COD} & \textbf{DIRECCIÓN ADMINISTRATIVA} & \textbf{MONTO (Bs)} & \textbf{\%} \\ 
\midrule
  \rowcolor{ctabla}
\DTLforeach{presupuesto}{%
    \vCod=COD, \vDesc=Descripcion, \vMonto=Monto, \vPorc=Porcentaje%
}{%
  \vCod & \vDesc & \vMonto & \vPorc\% \\\rowcolor{ctabla}

}
\textbf{--} & \textbf{TOTAL GENERAL EN BOLIVIANOS} & \textbf{1.737.058.874} & \textbf{100.00\%} \\
\end{tabular}
\end{small}


\section*{Fondo de atención Ciudadana (FAC 2025)}

\begin{small}
\renewcommand{\arraystretch}{0.8} % Reduce el espacio entre filas para que quepa en una página
\centering
\begin{tabular}{p{4cm} p{6cm} r r} 
\toprule
\textbf{Macrodistrito} & \textbf{Distrito} & \textbf{Monto Bs.} & \textbf{\%} \\ 
\midrule
\rowcolor{ctabla}


\DTLforeach{fac}{%
    \vMac=Macrodistrito, \vDist=Distrito, \vMonto=Monto, \vPorc=Porce%
}{%
    \vMac & \vDist & \vMonto & \vPorc\% \\\rowcolor{ctabla}

}

\textbf{--} & \textbf{TOTAL EN BOLIVIANOS} & \textbf{160.000.000} & \textbf{100.00\%} \\

\end{tabular}
\end{small}


\end{document}
